% ===== PRÉAMBULE CANONIQUE — à coller VERBATIM en tête de CHAQUE .tex =====
\documentclass[11pt,a4paper]{article}
\usepackage[utf8]{inputenc}\usepackage[T1]{fontenc}\usepackage{lmodern}
\usepackage[french]{babel}
\usepackage{amsmath,amssymb,mathtools}\usepackage{icomma}
\usepackage{siunitx}\sisetup{locale=FR}  % \qty{}{}, \si{}, \ang{}
\usepackage{xcolor}\usepackage{colortbl}
\usepackage{tikz}\usepackage{pgfplots}\pgfplotsset{compat=1.18}
\usepackage[siunitx,europeanresistors]{circuitikz} % circuits (élec.)
\usepackage[most]{tcolorbox}\usepackage{enumitem}\usepackage{array,booktabs}
\usepackage[a4paper,margin=2.2cm]{geometry}
\usetikzlibrary{arrows.meta,calc,angles,quotes,positioning,patterns,%
  backgrounds,shapes.geometric,decorations.pathmorphing,decorations.markings,intersections}
\definecolor{primary}{RGB}{222,49,99}\definecolor{primarydark}{RGB}{150,22,55}
\definecolor{defcol}{RGB}{0,114,178}\definecolor{methcol}{RGB}{52,92,148}
\definecolor{excol}{RGB}{0,135,90}\definecolor{attncol}{RGB}{200,40,55}
\tcbset{boxbase/.style={enhanced,breakable,boxrule=0.9pt,arc=2.5pt,
  left=3mm,right=3mm,top=2.3mm,bottom=2.3mm,fonttitle=\bfseries,coltitle=white}}
% --- boîtes TUTORIEL (noms EXACTS, ne pas renommer) ---
\newtcolorbox{cle}[1][]{boxbase,colback=defcol!6,colframe=defcol,
  title={\textsc{À retenir}\ifx\relax#1\relax\relax\else\textmd{ : \itshape #1}\fi}}
\newtcolorbox{methode}[1][]{boxbase,colback=methcol!7,colframe=methcol,sharp corners,
  title={\textsc{Méthode}\ifx\relax#1\relax\relax\else\textmd{ --- \itshape #1}\fi}}
\newtcolorbox{exemplebox}[1][]{boxbase,colback=excol!5,colframe=excol,
  title={\textsc{Exemple}\ifx\relax#1\relax\relax\else\textmd{ : \itshape #1}\fi}}
\newtcolorbox{piege}[1][]{boxbase,colback=attncol!6,colframe=attncol,
  title={\textsc{Piège}\ifx\relax#1\relax\relax\else\textmd{ : \itshape #1}\fi}}
\newtcolorbox{remarque}[1][]{boxbase,colback=black!4,colframe=black!45,
  title={\textsc{Remarque}\ifx\relax#1\relax\relax\else\textmd{ : \itshape #1}\fi}}
% --- boîtes EXERCICE / CORRIGÉ (noms EXACTS, auto-comptées) ---
\newtcolorbox[auto counter]{exo}[1][]{boxbase,colback=primary!4,colframe=primary,
  title={\textsc{Exercice}~\thetcbcounter\ifx\relax#1\relax\relax\else\textmd{ --- \itshape #1}\fi}}
\newtcolorbox[auto counter]{corrige}[1][]{boxbase,colback=excol!4,colframe=excol,
  title={\textsc{Corrigé}~\thetcbcounter\ifx\relax#1\relax\relax\else\textmd{ --- \itshape #1}\fi}}
\newcommand{\vv}[1]{\vec{#1}}\newcommand{\ds}{\displaystyle}
\newcommand{\R}{\mathbb{R}}\newcommand{\N}{\mathbb{N}}\newcommand{\Z}{\mathbb{Z}}
% (un \usepackage{fancyhdr}+en-tête est possible ; il est ignoré côté web)
% =========================================================================
\begin{document}
\sisetup{per-mode=symbol}

\begin{corrige}[deux charges de même signe]
\begin{enumerate}[label=\alph*)]
  \item \textbf{Même signe} (toutes deux positives ici). Aucune ligne ne va de l'une à l'autre : les
        lignes \emph{partent} des deux charges et s'écartent, ce qui n'arrive qu'entre deux charges
        de même signe (deux charges opposées seraient reliées par les lignes).
  \item Par symétrie, le point neutre est au \textbf{milieu} du segment $[AB]$ : les deux champs y
        sont opposés et de même intensité, donc $\vv{E} = \vv{0}$.
\end{enumerate}
\end{corrige}

\begin{corrige}[comparer deux charges par la densité]
\begin{enumerate}[label=\alph*)]
  \item La charge \textbf{$A$} : elle émet davantage de lignes (spectre plus dense), donc son champ
        est plus intense à distance égale.
  \item Le nombre de lignes est proportionnel à la charge :
        \[ \frac{q_A}{q_B} = \frac{8}{4} = 2. \]
        La charge $A$ est \textbf{deux fois} plus grande que $B$.
\end{enumerate}
\end{corrige}

\begin{corrige}[champ dans un condensateur]
\begin{enumerate}[label=\alph*)]
  \item \[ E = \frac{|U_{AB}|}{d} = \frac{200}{0{,}05} = \SI{4000}{\volt\per\metre}. \]
  \item $V_A > V_B$ : le champ est orienté de l'armature $A$ (potentiel élevé) vers l'armature $B$
        (potentiel le plus bas).
\end{enumerate}
\end{corrige}

\begin{corrige}[une charge lâchée dans le champ]
La force est $\vv{F} = q\,\vv{E}$, avec $\vv{E}$ vers le bas.
\begin{enumerate}[label=\alph*)]
  \item Charge \textbf{positive} : $\vv{F}$ a le \emph{même} sens que $\vv{E}$, elle part
        \textbf{vers le bas} (elle suit la ligne de champ).
  \item Charge \textbf{négative} : $\vv{F}$ est \emph{opposée} à $\vv{E}$, elle part \textbf{vers le
        haut} (à contre-sens de la ligne).
\end{enumerate}
\end{corrige}

\begin{corrige}[uniforme ou radial ?]
\begin{enumerate}[label=\alph*)]
  \item Le \textbf{spectre I} : ses lignes sont \emph{parallèles et équidistantes}, signature d'un
        champ uniforme. Le spectre II (lignes radiales qui s'écartent) est celui d'une charge
        ponctuelle, dont le champ \emph{décroît} avec la distance.
  \item Dans le champ uniforme, l'intensité $E$ est la \textbf{même en tout point} (même direction,
        même sens, même valeur) : les lignes restent également espacées partout.
\end{enumerate}
\end{corrige}

\end{document}
